\section{Prototipo de fidelidad media}
\label{sec:prototipo}

Esta sección presenta el prototipo de fidelidad media del simulador,
elaborado durante la PC02 como simulación visual del flujo principal
del producto. El prototipo cubre las siete pantallas que estructuran
el ciclo de uso del candidato y materializa las decisiones de diseño
documentadas en las secciones anteriores: el aura procedural en lugar
de paneles laterales (sección \ref{sec:alternativas}, dimensión 1),
el avatar 3D estilizado del entrevistador (dimensión 2), el panel
del observador en el modo \textit{peer mock} (dimensión 4) y la
gamificación basada en mejora relativa (módulo opcional 3 de la
sección \ref{sec:modulos}). Cada pantalla se enlaza con los
requerimientos funcionales y con los casos de uso a los que da
soporte.

\subsection{Pantalla de inicio}

La pantalla de inicio (figura \ref{fig:inicio}) es el primer punto de
contacto del usuario con la aplicación. Concentra el acceso a las
modalidades principales (entrevista individual y \textit{peer mock}),
muestra de forma resumida el estado del candidato (rango actual,
\textit{badges} recientes y próxima sesión sugerida) y mantiene la
identidad visual del \textit{Warachikuy} mediante una composición
centrada y serena. Cubre la entrada al caso de uso \cu{01} y la
navegación a \cu{04} (consultar historial) y \cu{05} (gamificación).

\begin{figure}[H]
\centering
\includegraphics[width=0.92\linewidth]{pantalla-inicio.png}
\caption{Pantalla de inicio del simulador. Punto de entrada a las
modalidades de entrevista y a las secciones de progreso y comunidad.}
\label{fig:inicio}
\end{figure}

\subsection{Pantalla de configuración inicial}

La configuración inicial (figura \ref{fig:configuracion}) recoge los
parámetros que condicionan la sesión: industria, rol objetivo, nivel
declarado y preferencias de accesibilidad (modo de estimulación
reducida, modo \textit{audio-first}, subtítulos, tipografía y
paciencia del entrevistador). El diseño separa el perfil profesional
del bloque de accesibilidad para que el segundo no quede como una
opción secundaria, en línea con el tratamiento de la accesibilidad
como requerimiento de primer nivel (\rnf{08}). Cubre el caso de uso
\cu{02} y los requerimientos \rf{17}--\rf{20}.

\begin{figure}[H]
\centering
\includegraphics[width=0.92\linewidth]{pantalla-configuracion-inicial.png}
\caption{Pantalla de configuración inicial. Reúne perfil profesional
y preferencias de accesibilidad antes de la primera sesión.}
\label{fig:configuracion}
\end{figure}

\subsection{Sala virtual --- vista del candidato}

La sala virtual desde la perspectiva del candidato
(figura \ref{fig:sala-candidato}) materializa el módulo de interfaz
gráfica no convencional: el avatar 3D estilizado del entrevistador
ocupa el centro de la composición y está rodeado por el aura
luminosa procedural que codifica las métricas multimodales
(\rf{05}). En lugar de \textit{dashboards} laterales, los elementos
de control y estado (privacidad, línea de tiempo circular, forma de
onda, miniatura del candidato y subtítulos opcionales) se reparten
discretamente por las esquinas. Cubre el flujo principal de \cu{01}
y los requerimientos \rf{02}, \rf{05}, \rf{06}, \rf{15}.

\begin{figure}[H]
\centering
\includegraphics[width=0.92\linewidth]{pantalla-sala-virtual-candidato.png}
\caption{Sala virtual --- vista del candidato. El aura procedural
alrededor del avatar reemplaza al \textit{dashboard} lateral
tradicional.}
\label{fig:sala-candidato}
\end{figure}

\subsection{Sala virtual --- vista del observador}

La vista del observador (figura \ref{fig:sala-publico}) corresponde
al modo \textit{peer mock} y cubre la decisión de la dimensión 4 de
la sección \ref{sec:alternativas}. El observador percibe en tiempo
real las métricas del candidato a través de una representación
simplificada del aura, dispone de un panel para anclar comentarios
en \textit{timestamps} específicos de la grabación y completa una
rúbrica estructurada por competencia al cierre de la sesión. La
distribución espacial conserva la del modo individual (heurística
H4 de la sección \ref{sec:usabilidad}: consistencia entre modos),
de modo que el cambio de rol no obliga a aprender una interfaz
nueva. Cubre \cu{03} y los requerimientos \rf{10}--\rf{12}.

\begin{figure}[H]
\centering
\includegraphics[width=0.92\linewidth]{pantalla-sala-virtual-publico.png}
\caption{Sala virtual --- vista del observador. Acceso instrumentado
a las métricas del candidato y captura de comentarios anclados.}
\label{fig:sala-publico}
\end{figure}

\subsection{Plan de mejora}

El plan de mejora (figura \ref{fig:plan-mejora}) se entrega al
cierre de cada sesión y consolida el resultado de la evaluación
multimodal. Presenta cuatro anillos circulares con el puntaje
agregado por competencia (comunicación verbal, lenguaje corporal,
contenido técnico, manejo del estrés), una lista priorizada de
ejercicios accionables generados por el LLM \textit{Coach} con
ejemplos concretos, y una línea de tiempo horizontal con los
momentos clave de la sesión que el candidato puede reproducir.
Cubre \cu{01} (paso 7) y \cu{04}, y los requerimientos \rf{07},
\rf{08} y \rf{09}.

\begin{figure}[H]
\centering
\includegraphics[width=0.92\linewidth]{pantalla-plan-mejora.png}
\caption{Pantalla del plan de mejora generado al cierre de la
sesión, con puntaje por competencia y ejercicios priorizados.}
\label{fig:plan-mejora}
\end{figure}

\subsection{Mi progreso}

La pantalla \textit{Mi progreso} (figura \ref{fig:mi-progreso})
materializa el módulo opcional de personalización continua. Reúne
la línea de tiempo de sesiones --donde cada nodo se pinta según la
valencia del aura final--, los \textit{sparklines} de tendencias por
competencia y el grid de \textit{badges} obtenidos. Un asistente
lateral permite al usuario interrogar al LLM sobre patrones del
historial (``muéstrame ejemplos para reducir muletillas'',
``compara mi última sesión con la mejor''). Cubre \cu{04} y los
requerimientos \rf{08}, \rf{09}, \rf{16}.

\begin{figure}[H]
\centering
\includegraphics[width=0.92\linewidth]{pantalla-mi-progreso.png}
\caption{Pantalla \textit{Mi progreso}. Visualización longitudinal
del desempeño del candidato y acceso al asistente lateral.}
\label{fig:mi-progreso}
\end{figure}

\subsection{Ranking y reto grupal}

La pantalla de ranking y reto grupal (figura \ref{fig:ranking})
materializa el módulo opcional de gamificación. La clasificación
semanal se calcula sobre la mejora relativa del usuario respecto a
su propia línea base, no sobre puntajes absolutos comparados entre
candidatos: la decisión evita penalizar a usuarios principiantes y
sostiene la motivación intrínseca, en línea con la heurística H8 de
diseño minimalista y con el atributo de satisfacción de la sección
\ref{sec:usabilidad}. La sección de retos grupales agrega objetivos
colectivos (acumulación compartida de horas o de \textit{peer mocks}
completados) con progreso visible para todos los miembros del grupo.
Cubre \cu{05} y los requerimientos \rf{13} y \rf{14}.

\begin{figure}[H]
\centering
\includegraphics[width=0.92\linewidth]{pantalla-ranking-reto-grupal.png}
\caption{Pantalla de ranking semanal por liga temática y sección de
retos grupales con progreso colectivo.}
\label{fig:ranking}
\end{figure}

\subsection{Diagrama de interacciones entre pantallas}

Las siete pantallas se conectan en el flujo descrito en la
figura~\ref{fig:flujo}. El usuario nuevo entra por la pantalla de
inicio, es derivado a la configuración inicial cuando aún no tiene
perfil, y desde allí accede al ciclo central
inicio~$\rightarrow$~sala virtual~$\rightarrow$~plan de mejora.
La sala virtual ofrece la bifurcación entre el modo individual
(vista del candidato) y el modo \textit{peer mock} (vistas del
candidato y del observador). El plan de mejora alimenta de forma
asíncrona a las pantallas \textit{Mi progreso} y \textit{Ranking},
que forman el ciclo longitudinal de personalización y gamificación.

\begin{figure}[H]
\centering
\includegraphics[width=0.95\linewidth]{diagrama-interacciones.jpeg}
\caption{Diagrama esquemático de las interacciones entre las
siete pantallas del prototipo.}
\label{fig:flujo}
\end{figure}

\subsection{Trazabilidad pantalla a caso de uso}

La tabla \ref{tab:trazabilidad-prototipo} consolida la cobertura del
prototipo respecto a los casos de uso primarios y a los
requerimientos funcionales clave. Las pantallas que no se prototipan
en esta entrega (administración del sistema, gestión de privacidad
y reporte institucional) corresponden a \cu{06}, \cu{07} y \cu{08},
y se incorporarán en iteraciones posteriores junto con los flujos
de moderación y exportación de datos.

\begin{table}[ht]
\centering\footnotesize
\renewcommand{\arraystretch}{1.25}
\begin{tabularx}{\linewidth}{@{}p{4.4cm} Y c@{}}
\toprule
\textbf{Pantalla} & \textbf{Requerimientos cubiertos} &
\textbf{Casos de uso} \\
\midrule
Inicio & \rf{01}, \rf{16} & \cu{01}, \cu{04}, \cu{05} \\
\addlinespace
Configuración inicial & \rf{17}--\rf{20}, \rnf{08} & \cu{02} \\
\addlinespace
Sala virtual --- candidato & \rf{02}, \rf{05}, \rf{06}, \rf{15} &
\cu{01} \\
\addlinespace
Sala virtual --- observador & \rf{10}, \rf{11}, \rf{12} & \cu{03} \\
\addlinespace
Plan de mejora & \rf{07}, \rf{08}, \rf{09} & \cu{01}, \cu{04} \\
\addlinespace
Mi progreso & \rf{08}, \rf{09}, \rf{16} & \cu{04} \\
\addlinespace
Ranking y reto grupal & \rf{13}, \rf{14} & \cu{05} \\
\bottomrule
\end{tabularx}
\caption{Trazabilidad de cada pantalla del prototipo a los
requerimientos funcionales y casos de uso correspondientes.}
\label{tab:trazabilidad-prototipo}
\end{table}

\subsection{Estado del prototipo y simulación del flujo}

El prototipo presentado en esta sección se entrega en formato de
fidelidad media: la composición espacial, la jerarquía visual y los
componentes principales (avatar, aura, línea de tiempo circular,
panel del observador, anillos de competencia) están definidos a un
nivel suficiente para validar las decisiones de diseño con usuarios
representantes de los perfiles primarios. La simulación clickeable
del flujo principal (\cu{01}) y el video corto del recorrido extremo
a extremo se publican junto al prototipo en el repositorio del
proyecto, en el directorio \texttt{docs/prototype/} del repo
público de GitHub. La iteración siguiente cubrirá las pantallas de
gestión de privacidad (\cu{06}), administración (\cu{07}) y reportes
institucionales (\cu{08}), que en esta PC02 se documentan en sus
casos de uso pero no se prototipan visualmente.
